In this paper, we presented \PipelineName, a systematic framework that reimagines medical video benchmark construction not as a labor-intensive manual task, but as a scalable, logic-driven engineering process. By bridging the gap between pixel-level visual extraction and symbolic knowledge representation, our pipeline successfully mitigates the hallucination risks inherent in purely generative approaches. 
% 本文提出了一种名为\PipelineName的系统化框架，它将医学视频基准测试的构建重新定义为一个可扩展的、逻辑驱动的工程流程，而非劳动密集型的手动任务。通过弥合像素级视觉提取和符号知识表示之间的鸿沟，我们的流程成功地降低了纯生成式方法中固有的幻觉风险。
The resulting dataset, \BenchmarkName, serves as a rigorous testbed for next-generation MLLMs, offering precise spatiotemporal grounding and verifiable multi-hop reasoning challenges that were previously attainable only through expert curation.
% 由此产生的数据集 \BenchmarkName 可作为下一代 MLLM 的严格测试平台，提供精确的时空基础和可验证的多跳推理挑战，而这些挑战以前只能通过专家整理才能实现。
Our empirical analysis confirms that the logical structures encoded by our pipeline effectively translate into the cognitive reasoning chains of SOTA models, proving that synthetic data can achieve both high fidelity and deep logic.
% 我们的实证分析证实，我们流程编码的逻辑结构有效地转化为最先进模型的认知推理链，证明合成数据可以同时实现高保真度和深度逻辑。